\section*{Capítulo \ref{ch:g7:literal} --- Expresiones literales}

\begin{solution}{exo:g7:literal:1}
$7x$; \quad $3y$; \quad $5ab$; \quad $x^2$; \quad $2(a + 4)$.
\end{solution}

\begin{solution}{exo:g7:literal:2}
$5ab = 5 \times a \times b$; \quad
$3(x+2) = 3 \times (x + 2)$; \quad
$x^2 + 4x = x \times x + 4 \times x$; \quad
$2a^3 = 2 \times a \times a \times a$.
\end{solution}

\begin{solution}{exo:g7:literal:3}
Para $x = 3$: \ $5 \times 3 + 1 = 16$; \quad $3 \times 3 = 9$; \quad
$2 \times 9 - 3 = 15$; \quad $10 - 6 = 4$.
\end{solution}

\begin{solution}{exo:g7:literal:4}
$3a + 2b = 12 + 3 = 15$; \qquad
$ab + 5 = 4 \times 1.5 + 5 = 6 + 5 = 11$.
\end{solution}

\begin{solution}{exo:g7:literal:5}
$8x + 5x = 13x$; \quad
$9a - 4a + a = 6a$; \quad
$6y + 2 + 3y + 7 = 9y + 9$; \quad
$5t + 3 - 2t - 3 = 3t$.
\end{solution}

\begin{solution}{exo:g7:literal:6}
Cuadrado: $P = 4c$. \omterm{def:g6:shapes:triangles}{Triángulo isósceles}: $P = b + 2s$. Panes:
$p = 1.20\,n$ euros.
\end{solution}

\begin{solution}{exo:g7:literal:7}
$2n$: el doble; \ $n + 2$: dos más que $n$; \ $n^2$: el cuadrado;
\ $\frac n2$: la \omterm{def:g3:division:half}{mitad}.
\end{solution}

\begin{solution}{exo:g7:literal:8}
$3(x+4) = 3x + 12$: \emph{verdadera} para todo $x$ (distributividad ---
esto es una demostración, no hace falta probar valores).

$2x + 5 = 7x$: prueba con $x = 2$: izquierda $9$, derecha $14$ --- \emph{falsa} en
general (solo se cumple para $x = 1$).

$x + x = x^2$: prueba con $x = 3$: izquierda $6$, derecha $9$ --- \emph{falsa} en
general ($x + x = 2x$).
\end{solution}

\begin{solution}{exo:g7:literal:9}
\emph{1.} Largo: $a + 3$. \omterm{def:g6:measure:perimeter}{Perímetro}:
\[
P = 2\bigl(w + (w + 3)\bigr) = 2(2w + 3) = 4w + 6 .
\]

\emph{2.} Para $a = 5$: directamente, los lados son $5$ y $8$, así que
$P = 2 \times 13 = 26$ cm; con la fórmula, $4 \times 5 + 6 = 26$ cm.
Coinciden.
\end{solution}

\begin{solution}{exo:g7:literal:10}
Siguiendo los pasos con la letra $n$:
\[
n \ \to\ n + 6 \ \to\ 3(n + 6) = 3n + 18 \ \to\ 3n + 18 - 18 = 3n
\ \to\ \frac{3n}{n} = 3 .
\]
Todo el mundo obtiene $3$.
\end{solution}

\begin{solution}{exo:g7:literal:11}
Como \omterm{ex:g1:subtraction:difference}{diferencia}: $A = a^2 - b^2$. Cortando la L en dos
rectángulos: una franja de ancho completo y altura $a - b$ y, encima, un
rectángulo más estrecho:
\[
A = a(a - b) + b(a - b)
\]
(el rectángulo de abajo, de ancho $a$ y alto $a - b$, y el de arriba, de ancho $a - b$
y alto $b$, dan la variante $a(a-b) + (a-b)b$ --- lo mismo).
Para $a = 5$ y $b = 2$: $a^2 - b^2 = 25 - 4 = 21$, y
$5 \times 3 + 2 \times 3 = 15 + 6 = 21$. Las dos expresiones coinciden ---
y siempre lo hacen: es la identidad $a^2 - b^2 = (a + b)(a - b)$ del
\cref{ch:g9:algebra} disfrazada.
\end{solution}

\begin{solution}{pb:g7:literal:1}
\textbf{1.} Las figuras $1$ a $4$ necesitan $4$, $7$, $10$ y $13$
cerillas.

\textbf{2.} Un cuadrado nuevo comparte un lado con la figura
anterior, así que solo $3$ de sus $4$ lados son cerillas nuevas. Desde la
figura $4$ ($13$ cerillas), seis cuadrados más cuestan
$6 \times 3 = 18$: la figura $10$ necesita $13 + 18 = 31$ cerillas.

\textbf{3.} Leyendo la fila desde la izquierda: empieza con una
cerilla vertical y después cada uno de los $n$ cuadrados aporta un
lado de arriba, uno de abajo y uno de la derecha --- $3$ cerillas cada uno. Total:
$3n + 1$; el «$+1$» es la primerísima cerilla vertical. Comprobación:
$n = 1, 2, 3, 4$ dan $4, 7, 10, 13$.

\textbf{4.} $4 + 3(n - 1) = 4 + 3n - 3 = 3n + 1$: la fórmula de Amir
se desarrolla exactamente en la misma expresión reducida. Coincidencia
total.

\textbf{5.} Lena: $2n + (n + 1) = 3n + 1$ --- otra vez lo mismo.
Tres maneras de verlo y una sola forma reducida. Figura $100$:
$3 \times 100 + 1 = 301$ cerillas.

\textbf{6.} Recuentos: $3$, $5$ y $7$ cerillas. Cada triángulo nuevo
se apoya en un lado ya existente y solo añade $2$ cerillas, partiendo
de $3$: la figura $n$ necesita $3 + 2(n - 1) = 2n + 1$ cerillas.

\textbf{7.} Deshaciendo $2n + 1 = 49$: quitando la primera cerilla quedan
$48$, y cada triángulo se lleva $2$:
$48 \div 2 = 24$ triángulos.

\textbf{8.} En una fila de $n$ mesas, cada mesa ofrece su lado
de arriba y el de abajo ($2n$ plazas), y las dos mesas de los extremos ofrecen un
lado más cada una ($+2$): $2n + 2$ plazas. Comprobación: $4$, $6$ y $8$
para $n = 1, 2, 3$.

\textbf{9.} Una sola fila: $2n + 2 \geq 30$ exige $2n \geq 28$, así que
$14$ mesas. Repartiendo esas mismas $14$ mesas en dos filas de
$7$: cada fila sienta a $2 \times 7 + 2 = 16$ y, juntas, a $32$ ---
dos más que la fila única, porque cada fila nueva aporta dos
\emph{extremos} nuevos. (Cada división gana $2$ plazas; los cáterin lo saben.)

\textbf{10.} La L de tamaño $n$ tiene $n$ puntos hacia arriba y $n$
hacia el lado, compartiendo el punto de la esquina: $n + n - 1 = 2n - 1$ puntos.
Figuras $1$, $2$ y $3$: $1$, $3$ y $5$ puntos --- los \emph{\omterm{def:g3:numbers:evenodd}{números
impares}}. Encajadas una alrededor de la siguiente, las L llenan un
cuadrado de $n \times n$ puntos: la historia del
\cref{pb:g8:equations:1}.

\textbf{11.} $3(n + 2) - 5 = 3n + 6 - 5 = 3n + 1$: la afirmación
se reduce a la fórmula demostrada, así que es correcta para todo $n$ ---
establecida con álgebra, no con ejemplos.

\textbf{12.} Pruebas: $4 \geq 2$, $9 \geq 3$, $100 \geq 10$: todas
pasan. Pero $n = 0.5$ da $n^2 = 0.25 < 0.5$: la afirmación falla.
Las tres pruebas no demostraron nada sobre \emph{todos} los números (solo
no lograron encontrar problemas); el único contraejemplo
\emph{demuestra} que la afirmación general es falsa
(\cref{met:g7:literal:test}). Una letra representa a todos los números
--- decimales incluidos.

\textbf{13.} Regiones: $2$ puntos $\to 2$; $3$ puntos $\to 4$;
$4$ puntos $\to 8$; $5$ puntos $\to 16$. El patrón susurra
«cada punto duplica el recuento»: se esperarían $32$ regiones para
$6$ puntos.

\textbf{14.} Un recuento cuidadoso da $31$ regiones --- no $32$.
La conjetura del doble ha muerto: sobrevivió a cuatro comprobaciones y
falló en la quinta. (Este contraejemplo es tan famoso que tiene
nombre, «la trampa de las regiones del círculo»; nada de atajos con puntos regulares
--- tres cuerdas que se cortasen en un mismo punto fundirían regiones y
estropearían el recuento.)

\textbf{15.} Una demostración habría necesitado una razón \emph{estructural}
de por qué añadir un punto duplica las regiones --- y no la
hay: la duplicación era una casualidad de los casos pequeños. Las
fórmulas de las cerillas eran seguras porque cada una se leía en la
propia construcción (cada cuadrado cuesta visiblemente $3$ cerillas y
cada mesa ofrece visiblemente $2$ plazas): estructura, no pruebas.
Predicción demostrada: $3 \times 2\,026 + 1 = 6\,079$ cerillas.
\end{solution}
